\documentclass[12pt, twoside, openright]{report} % Fuente a 12pt, formato doble página y chapter a la derecha
\raggedbottom % No ajustar el contenido con un salto de página

% MÁRGENES: 2,5 cm sup. e inf.; 3 cm izdo. y dcho.
\usepackage[
a4paper,
vmargin=2.5cm,
hmargin=3cm
]{geometry}

% INTERLINEADO: Estrecho (6 ptos./interlineado 1,15) o Moderado (6 ptos./interlineado 1,5)
\renewcommand{\baselinestretch}{1.15}
\parskip=6pt

% DEFINICIÓN DE COLORES para portada y listados de código
\usepackage[table]{xcolor}
\definecolor{azulUC3M}{RGB}{0,0,102}
\definecolor{gray97}{gray}{.97}
\definecolor{gray75}{gray}{.75}
\definecolor{gray45}{gray}{.45}

% Soporte para GENERAR PDF/A
\usepackage{etoolbox}
\makeatletter
\@ifl@t@r\fmtversion{2021-06-01}%
 {\AddToHook{package/after/xmpincl}
   {\patchcmd\mcs@xmpincl@patchFile{\if\par}{\ifx\par}{}{\fail}}}{}
\makeatother
\usepackage[a-1b]{pdfx}

% ENLACES
\usepackage{hyperref}
\hypersetup{colorlinks=true,
  linkcolor=black, % enlaces a partes del documento (p.e. índice) en color negro
  urlcolor=blue} % enlaces a recursos fuera del documento en azul

% Añadir pdfs como partes del documento
\usepackage{pdfpages}

% Quitar la indentación de principio de los párrafos
\setlength{\parindent}{0em}

% EXPRESIONES MATEMÁTICAS
\usepackage{amsmath,amssymb,amsfonts,amsthm}

\usepackage{txfonts} 
\usepackage[T1]{fontenc}
\usepackage[utf8]{inputenc}

% Insertar gráficas y fotos
\usepackage{tikz}
\usetikzlibrary{cd}
\usepackage{pgfplots}

\usepackage[spanish, es-tabla]{babel} 
\usepackage[babel, spanish=spanish]{csquotes}
\AtBeginEnvironment{quote}{\small}

% diseño de PIE DE PÁGINA
\usepackage{fancyhdr}
\pagestyle{fancy}
\fancyhf{}
\renewcommand{\headrulewidth}{0pt}
\fancyfoot[LE,RO]{\thepage}
\fancypagestyle{plain}{\pagestyle{fancy}}

% DISEÑO DE LOS TÍTULOS de las partes del trabajo (capítulos y epígrafes o subcapítulos)
\usepackage{titlesec}
\usepackage{titletoc}
\titleformat{\chapter}[block]
{\large\bfseries\filcenter}
{\thechapter.}
{5pt}
{\MakeUppercase}
{}
\titlespacing{\chapter}{0pt}{0pt}{*3}
\titlecontents{chapter}
[0pt]                                               
{}
{\contentsmargin{0pt}\thecontentslabel.\enspace\uppercase}
{\contentsmargin{0pt}\uppercase}                        
{\titlerule*[.7pc]{.}\contentspage}                 

\titleformat{\section}
{\bfseries}
{\thesection.}
{5pt}
{}
\titlecontents{section}
[5pt]                                               
{}
{\contentsmargin{0pt}\thecontentslabel.\enspace}
{\contentsmargin{0pt}}
{\titlerule*[.7pc]{.}\contentspage}

\titleformat{\subsection}
{\normalsize\bfseries}
{\thesubsection.}
{5pt}
{}
\titlecontents{subsection}
[10pt]                                               
{}
{\contentsmargin{0pt}                          
  \thecontentslabel.\enspace}
{\contentsmargin{0pt}}                        
{\titlerule*[.7pc]{.}\contentspage}  

\usepackage{adjustbox}

% DISEÑO DE TABLAS.
\usepackage{multirow} % permite combinar celdas 
\usepackage{caption} % para personalizar el título de tablas y figuras
\usepackage{floatrow} % utilizamos este paquete y sus macros \ttabbox y \ffigbox para alinear los nombres de tablas y figuras de acuerdo con el estilo definido. Para su uso ver archivo de ejemplo 
\usepackage{array} % con este paquete podemos definir en la siguiente línea un nuevo tipo de columna para tablas: ancho personalizado y contenido centrado
\newcolumntype{P}[1]{>{\centering\arraybackslash}p{#1}}
\DeclareCaptionFormat{upper}{#1#2\uppercase{#3}\par}

% Diseño de tabla para ingeniería
\captionsetup[table]{
  format=hang,
  name=Tabla,
  justification=centering,
  labelsep=colon,
  width=.75\linewidth,
  labelfont=small,
  font=small,
}

% DISEÑO DE FIGURAS.
\usepackage{graphicx}
\graphicspath{{img/}} %ruta a la carpeta de imágenes

% Diseño de figuras para ingeniería
\captionsetup[figure]{
  format=hang,
  name=Fig.,
  singlelinecheck=off,
  labelsep=colon,
  labelfont=small,
  font=small    
}

% NOTAS A PIE DE PÁGINA
\usepackage{chngcntr} % Para numeración continua de las notas al pie
\counterwithout{footnote}{chapter}

% LISTADOS DE CÓDIGO
% soporte y estilo para listados de código. Más información en https://es.wikibooks.org/wiki/Manual_de_LaTeX/Listados_de_código/Listados_con_listings
\usepackage{listings}

% definimos un estilo de listings
\lstdefinestyle{estilo}{ frame=Ltb,
  framerule=0pt,
  aboveskip=0.5cm,
  framextopmargin=3pt,
  framexbottommargin=3pt,
  framexleftmargin=0.4cm,
  framesep=0pt,
  rulesep=.4pt,
  backgroundcolor=\color{gray97},
  rulesepcolor=\color{black},
  %
  basicstyle=\ttfamily\footnotesize,
  keywordstyle=\bfseries,
  stringstyle=\ttfamily,
  showstringspaces = false,
  commentstyle=\color{gray45},     
  %
  numbers=left,
  numbersep=15pt,
  numberstyle=\tiny,
  numberfirstline = false,
  breaklines=true,
  xleftmargin=\parindent
}

\captionsetup[lstlisting]{font=small, labelsep=period}
% fijamos el estilo a utilizar 
\lstset{style=estilo}
\renewcommand{\lstlistingname}{\uppercase{Código}}

\pgfplotsset{compat=1.17} 
%-------------
% DOCUMENTO
%-------------

\begin{document}
\pagenumbering{roman} % Se utilizan cifras romanas en la numeración de las páginas previas al cuerpo del trabajo

%----------
% PORTADA
%---------- 
\begin{titlepage}
  \begin{sffamily}
    \color{azulUC3M}
    \begin{center}
      \begin{figure}[H] % Incluimos el logotipo de la Universidad
        \makebox[\textwidth][c]{\includegraphics[width=16cm]{Portada_Logo.png}}
      \end{figure}
      \vspace{2.5cm}
      \begin{Large}
        Grado en Ingeniería Informática\\
        2021-2022\\
        \vspace{2cm}
        \textsl{Apuntes}\\
        \bigskip
      \end{Large}
      {\Huge Mitos Griegos y Cultura Europea}\\
      \vspace*{0.5cm}
      \rule{10.5cm}{0.1mm}\\
      \vspace*{0.9cm}
      {\LARGE Jorge Rodríguez Fraile\footnote{\href{mailto:100405951@alumnos.uc3m.es}{Universidad: 100405951@alumnos.uc3m.es}  |  \href{mailto:jrf1616@gmail.com}{Personal: jrf1616@gmail.com}}}\\
      \vspace*{1cm}
    \end{center}
    \vfill
    \color{black}
    \includegraphics[width=4.2cm]{img/creativecommons.png}\\
    Esta obra se encuentra sujeta a la licencia Creative Commons\\ \textbf{Reconocimiento - No Comercial - Sin Obra Derivada}
  \end{sffamily}
\end{titlepage}

%----------
% ÍNDICES
%---------- 

%--
% Índice general
%-
\tableofcontents
\thispagestyle{fancy}

%--
% Índice de figuras. Si no se incluyen, comenta las líneas siguientes
%-
\listoffigures
\thispagestyle{fancy}

%--
% Índice de tablas. Si no se incluyen, comenta las líneas siguientes
%-
\listoftables
\thispagestyle{fancy}

%----------
% TRABAJO
%---------- 

\pagenumbering{arabic} % numeración con números arábigos para el resto de la publicación  

%----------
% COMENZAR A ESCRIBIR AQUÍ
%---------- 

\chapter{Información}\label{ch:informacion}
\section{Profesores}\label{sec:profesores}
\begin{quote}
	Magistral: Lorena Pérez Yarza, loperezy@inst.uc3m.es
\end{quote}

\section{Sistema de evaluación}\label{sec:sistema-de-evaluación}
Hay evaluación continua, que si se sigue no hace falta ir a un examen final.
\begin{enumerate}
    \item Visita del Prado (17 de febrero) y Práctica
    \begin{itemize}
      \item Visita corta con la Uni (jueves por la mañana/sábado por la mañana)
      \item Hacer un trabajo de 2-3 páginas (3 páginas máximos), con imágenes, explicando 3 obras de arte con contenido mitológico.
    \end{itemize}
  \item Seminario online: 17 de febrero o 3 de marzo (es una alternativa al museo).
  \item 1 lectura obligatoria: Selección de obras en aula global.
    \begin{itemize}
      \item Extensión de 3 páginas (máximo 3 páginas), pero no solo copiar y pegar.
      \item Recensión: Decir en que consiste la obra y que impresión nos ha causado, como reflexión personal, teniendo en cuenta el ámbito mitológico.
      \item Decir más o menos como se organiza, cuál es la temática y como aborda esa temática (irónica, burlona, poética, etc.). Explicar el vehículo, el libro, el cómo se explican las cosas es lo que interesa de una recensión. Cuanto es de carácter científico se evalúa si es bueno o malo el libro según si se adecua a la cuestión, expresando que entendemos el texto y que sabemos explicarlo de forma resumida, pero sin incluir todos los hechos, hacer una selección de ejemplos importantes para que quede racional.
    \end{itemize}
  \item (Examen final) Presentaciones finales: oral en grupo y trabajo individual (18 de mayo, límite)
    \begin{itemize}
      \item Un tema mitológico que aparezca en la actualidad, ya sea libro, película, comic o contexto social.
      \item Sobre ese tema hay que explicar de que va el mito antiguo y su origen, para después mostrar en que ámbitos o lugares se da en la actualidad.
      \item Nos podemos inspirar en las Transformaciones de Ovidio, salen historias muy famosas que podemos usar.
      \item Grupos de 2 personas máximo para la presentación.
      \item Presentación oral de 5 minutos, online.
      \item Trabajo individual de 5-8 páginas (5 páginas mínimo) sobre este tema, entre compañeros será parecido. Además, debe incluir una reflexión personal tanto del mito como de su uso en la actualidad (con ejemplos) y una bibliografía. Puede incluir imágenes.
    \end{itemize}
\end{enumerate}

\chapter{Introducción}\label{ch:introduccion}
\textbf{Lord Byron}, poeta romántico británico, combatió en la Guerra de independencia griega y murió en esta.

La mitología era un parte cultural oral y viva que se usaba para explicar el mundo que les rodea, servían para enseñar. Entender la mitología griega nos puede ayudar a entender el mundo antiguo y su historia, la historia de los museos, edificios históricos o incluso analizar arte del cine o los cuadros.

\textbf{Heinrich Schliemann (1822-1890)} se inspiró en los relatos míticos y participo en los relatos míticos y participó activamente en el descubrimiento arqueológico de Troya y Micenas, lugares de la Ilíada. Dando lugar a que se tuviera más en cuenta los mitos y que la historia no solo se reconstruye con hechos fácticos.

Los historiadores y filósofos consideraban que eran la pura negación de la historia, ya que no tiene verosimilitud y son muy fantásticos, carecen de una lógica formal. Se pensaba que la ciencia tenía que explicar todo, pero no es así las personas somos irracionales que actúan por impulsos y no todo tiene que ser fáctico.

En muchos casos se dan hechos que no son probados, pero que se han mitificado hasta el punto de que todo el mundo piense que es cierto. Es el caso del Cid o la Guerra de Troya.

En principio eran historias sagradas de la tribu que explicaban el mundo mediante relatos de hechos memorables donde dioses y héroes de antaño actuaban. Historias de otro tiempo, de los orígenes del mundo, separado del nuestro mentalmente que por acciones divinas y heroicas en una época primordial cambiaban, lo ordenaban o lo rediseñaban hasta ser el que tenemos hoy.

"Los mitos 'clásicos' representan ya el triunfo de la obra literaria sobre la creencia religiosa" \textbf{(Mircea Eliade, 1994: 166)}

Eran reescritos y modificados por escribas, dando lugar a veces a confusión y prejuicios de los historiadores, dando el significado que ellos piensan. Por ejemplo la eliminación de algunos dioses menores, como el Dios Sol fue eliminado durante el renacimiento por los escribas al pensar que era un dios oriental y no romano, que era una confusión y sustituyéndolo por Júpiter.

"Los mitos viven en el país de la memoria" \textbf{(Marcel Detienne)}

Se transmite de manera oral normalmente de generación en generación, aunque también se hace de manera literaria en instituciones especializadas. Cuando los griegos aprendían a escribir/leer lo primero que aprendían era la Ilíada, era conocido por todos.

Son ejes culturales, tanto políticamente como ligados a las creencias, tienen el mismo objetivo, van de la mano, explican ciertos ritos y acontecimientos.

La cultura occidental no es pagana, bebe de la memoria cultural romana y griega, cultura clásica, pero en muchas zonas se ha expandido y adaptado a las diferentes culturas. Por ejemplo los campos elíseos, como lugar bonito tras la muerte bebe de la cultura egipcia.

\textbf{Mitema (tópico literario):} En el estudio de la mitología es la parte más pequeña que conforma un mito (equivale al fonema de la lingüística), es un elemento constante que siempre aparece intercambiado y reensamblado con otros mitemas o contextos. Son partes estructurales. Puede ser la idea de un viaje de regreso por el mar o un tipo de nacimiento/hecho milagroso (ejem. echado a un río, sumerios con Sargón de Acad o cristianismo con Moisés), el amor, un diluvio (ejem. sumerios y cristianismo). Son buenas ideas, que sirven para articular la obra y que cuando pasan a estar por escrito son tópicos literarios.

El mito se transmite a partir de ideas exitosas, pueden ser mitemas o ideas de contexto generales que gustan, que se usan para formar nuevas palabras o seres fantásticos como quimera que significaba cabra y se interpreta como una bestia ceremonial. Por ejemplo capricornio parte del calendario. Pasa con las sirenas, grifos, etc.

Los mitos eran parte de la cultura viva e iban evolucionando constantemente, hay que tener muy presente que tienen una larga historia y variaciones. Tienen muchas raíces y es difícil llegar hasta el original, diferenciando obras literarias, teatrales o religiosas.

Hay muchos animales a los que se les dota de significado, como el Pegaso, Minotauro, Hidra, Paloma, León, Toro, Esfinge, etc.

El mito es fundamentalmente la historia oral, más que la escrita y literaria.

\chapter{Clase 2}\label{ch:clase2}
Un ejemplo de mito son Las Tres Gracias, que representan Aglaya (‘Belleza’), Eufrósine (‘Júbilo’) y Talia (‘Abundancia’), suelen ser acompañantes de Venus. Sus nombres representan el concepto que representan. Salen en el mito de Marte destruyendo todo incluidas las Gracias. Se han representado de muchas maneras a lo largo de la historia.

Es muy habitual la representación en tríos como las Moiras (hilanderas del destino) o las Horas (naturaleza y estaciones), que representan ciertos conceptos alegóricos y se relacionan con los dioses.

También están las musas que son 9 y representan seres divinos que inspiran al arte.

¿Dónde encontramos el mito?
\begin{enumerate}
  \item En la alegoría filosófica, desde la cosmogonía (el inicio del universo) hasta temas filosóficos.
  \item La explicación racional de los mitos (libros eruditos).
  \item En sátira o parodia.
  \item Fábula sentimental.
  \item En el teatro (Ejem. Edipo Rey o Antígona)
  \item Relatos históricos y geográficos, que indirectamente se refieren a los hechos de los mitos o en las propias religiones.
  \item Propaganda política, se apropian un dios o héroe que les represente.
  \item Sagas heroicas y epopeyas.
  \item Explicaciones complementarias en ficción realista, como La Ilíada.
  \item En moralejas antiguas y cuentos moralizantes, como La Odisea, emplean partes de los mitos para hacer una moraleja, tal como Ícaro con aspirar a algo demasiado alto.
  \item Anécdotas curiosas o humorísticas en obras variadas, como Luciano de Samosata que habla de los mitos con un tono cínico.
  \item En representaciones artísticas, las más amplias, al haber tantas nos permite saber cuáles eran las más populares y más detalles al haber tanto escrito.
\end{enumerate}

Ejem. de mitos en la actualidad: Diosa de la justicia en tribunales, maratón, juegos olímpicos, astronomía, zodiaco, logo Nike, nombres de los planetas (larga historia), continente Europa, etc.

En muchos casos se cuentan historias uniendo y relacionando los mitos, se reutilizan de manera constante por las diferentes culturas.

\textbf{Mitopoiesis} consiste la creación de un mundo fantásticos/ficticios basando en hechos o una referencia. Es el género narrativo en el cual el autor crea todo un conjunto de conceptos, regiones, personajes, sucesos, y arquetipos interrelacionados creando una mitología propia.

Sabemos de los mitos por los Aedos, cantantes populares que creaban canciones con ritos fáciles de memorizar con un bastón, y mitógrafos, redactores de mitos en épocas populares (poetas, filósofos) y recolectores de estos.

Hay autores que crearon sus propios mitos y otros se encargaban de crear un corpus común de las historias.

900-600 a.C. Periodo arcaico de Grecia, anterior a que se pusiera la historia por escrito.

Cubrían la zona de Grecia y Turquía, posteriormente se expandieron, algunos los utilizaban los mitos como inspiración y otros los escribían propiamente.

323 - 31 a.C. Época Helenística, la consolidación escrita de la mitología.

La primera epopeya que se tiene es la Gilgamesh, que se ha conseguido reproducir en la actualidad de la manera que se hacía en la época. Narraciones lentas para remarcar los hechos épicos, con mucha repetición para facilitar que la gente lo recuerde.

Olimpo, el hogar de los dioses.

Se pensaba que los dioses habitaban en lugares inaccesibles, ya sea en las montañas (el Olimpo), cielo, subterráneo, etc. De esta manera están en un lugar que los morales no alcanzaban, en la época no tenían los medios, es por esto que en los caminos había santuarios.

Estos lugares tenían un sentido simbólico, por ejemplo monte Olimpo había varios.

\textbf{Cosmogonía:} Narración mítica o un modelo, que pretende dar respuesta al origen del universo y de la propia humanidad.

\textbf{Teogonía:} Contiene una de las más antiguas versiones del origen del cosmos y el linaje de los dioses de la mitología griega. Habla de la creación del mundo.

Son temas muy recurrente por los aedos, en sus cantos.

Hay muchas versiones del origen del mundo.

Rea y Cronos son los que dan origen a los dioses, por una profecía de que su hijo acabara con él se los come. Rea arta oculta a Zeus en Creta. Un día se traga una piedra y vomita a todos los dioses dándoles origen.

La teogonía de Hesíodo:
\begin{enumerate}
  \item Mundo en contacto, Koiné fenicia (la puesta en común fenicia).
  \item Poeta rural de Beocia. \\ Musas.
  \item Competidor en certámenes de poesía.
  \item Cronología incierta.
  \item Recoge parte del conocimiento de su tiempo.
  \item En representaciones artísticas.
\end{enumerate}

El origen de Venus, se puede ver en El nacimiento de Venus de Botticelli 1482-1485, Afrodita nace de los genitales de Urano cortado por Cronos y al llegar a tierra.

Las grandes luchas: Titanomaquia
\begin{enumerate}
  \item Caída de la primera dinastía.
  \item Dioses primordiales (conceptuales)
  \item Cambios de dinastías de dioses.
  \item El episodio se reproduce con la Gigantomaquia y lucha contra Tifón.
\end{enumerate}

Se utilizaban los mitos para dar nombre a algunas cosas o ciertas cosas inspiraban para crear mitos.


\chapter{Clase 3}\label{ch:clase3}
El mito es atemporal, no sigue un hilo argumental fijo.

La relación entre los dioses/lo sobrenatural se relaciona con los mortales/lo natural.

Ganímedes, fue raptado por Zeus y se convierte en el copero de los dioses. Zeus usa sus atributos, como águila, para conseguir un nuevo efebo. En la época el amor y el matrimonio van por separado, muchas veces se representa con cadenas, que te obligan a hacer cosas que no quieres.

Cuando hay una relación entre dos hombres, se muestra una clara diferencia social, la sumisión es un elemento importante. Por ejemplo, Adriano que tras el ahogo de su amado construye estatuas y lo trata de divinizar, como paso con Ganímedes.

\section{Zeus - Júpiter}
\begin{itemize}
	\item Señor, rey y padre
	\item Señor del Cielo y tormenta
	\item Dios de la Justicia (con Metis = Prudencia)
	\item Zeus Olímpico, Horkios, Xenios\dots
	\item Faceta mitológica
	\item Faceta de dios supremo
	\item Evolución cronológica
	\item Sus símbolos son el águila, el trono, el rayo y el cetro (que sería usado por los reyes).
\end{itemize}

Zeus se come a su primera mujer, Metis, para que no le pase lo mismo que a su padre y abuelo, pero le provoca gran dolor y pedirá que le rompan la cabeza, donde nace Atenea.

Se les designan epítetos, conjuntos de adjetivos, a los dioses, reyes y héroes para definirlos, su grandeza. Además, sirven para mostrar respeto. Atributos como \enquote{el Grande}, \enquote{el Sabio}\dots

Se habla tanto de los amoríos como manera de tejer la genealogía y relación entre los dioses.

Uno de los santuarios más importantes es: el Santuario de Olimpia, donde tiene lugar las olimpiadas clásicas, se iba ampliando con los años para realizar más juegos. En el Templo de Zeus era permanente y se encontraron estatuillas del neolítico. Además, el monte más próximo es el Kronos. Se celebraron los juegos desde 700 a.C. hasta 300 d.C. que con la religión se detuvieron.

Tras la victoria de los dioses se dividen los poderes entre ellos, Hades el Inframundo (mundo subterráneo), Poseidón el señor de los océanos (, caballos y terremotos) y Zeus el dios del cielo. Y todos mandan sobre la tierra, donde viven los mortales.
El Panteón está formando por todos los dioses, 12 olímpicos y otros que varían según las escrituras.

\section{Más allá de Tifón}
Tras esto tenemos a las diosas del destino, seres divinos (no diosas) de poder primordial, las Moiras (Nona, Decima y Morta). Representan las diferentes fases de la vida.

Las Furias (nacen de la sangre de Urano), las Harpías y Seres variados salvajes (como Centauros).

Centauros, representan las pasiones humanas sin control; lujuriosos, violentos y borrachos. Quirón es un ejemplo de centauro civilizado, que entreno a grandes héroes y dioses para la batalla, muere de manera accidental con una flecha envenenada que le hiere eternamente (renunciando por su dolor a la inmortalidad).

\section{Prometeo}
\begin{itemize}
	\item Titán benévolo
	\item Dador de fuego (civilización)
	\item Por él, el sacrificio es como es
	\item Engañó con astucia a Zeus y roba el fuego a los dioses.
\end{itemize}

\section{La Caja de Pandora}
\begin{itemize}
	\item Uno de los mitos de origen del género humano
	      \begin{itemize}
		      \item Deucalión y primera
		      \item Barro
		      \item Pandora (significa \enquote{la que da todas las cosas} fue dotada por los dioses de sus virtudes
	      \end{itemize}
	\item Originalmente era una jarra y no una caja.
	\item Epimeteo, hermano de Prometeo al que le regalan a Pandora.
	\item Construcción de una mujer, es la primera mujer.
	\item Visión misógina
	\item Origen de los males, al estar dotada de curiosidad no se puede resistir a abrirla y libera a todos los males
\end{itemize}

Un mitema muy habitual es crear al hombre a partir del barro.

\section{El Castigo de Prometeo}
\begin{itemize}
	\item Encadenado a una montaña en Escitia (Cáucaso), el lugar más lejano. Donde un águila, el símbolo de Zeus el que le castiga, se comería sus vísceras cada día, como ser inmortal las regenera, un castigo brutal.
	\item Esquilo, Prometheia
	      \begin{itemize}
		      \item Prometero Encadenado
		      \item Prometero liberado
		      \item Prometeo portador del fuego
	      \end{itemize}
	\item Prometeo capturado, de Rubens
	\item Diferentes interpretaciones > diferentes
	      \begin{itemize}
		      \item Hesíodo
		      \item Platón
		      \item Esquilo
		      \item Pausanias: cultos
	      \end{itemize}
\end{itemize}

La liberación de Prometeo, Hércules/Heracles es el encargado de liberarle y vencer al águila.

Heracles es el héroe por antonomasia (buen ejemplo para trabajo final), pasa de héroe a dios, era hijo de Zeus y Almenada. Al ser hijo bastardo de Zeus, es mitad dios, mitad mortal. Destaca por su fortaleza y buen corazón, aunque no era muy listo y se metía en problemas.

Es la base para los héroes modernos, tiene que realizar proezas para compensar su poder.

Zeus engaña a Hera para que amante a Heracles, pero este al morder tan fuerte le hace daño y nota que va ganando fuerza e inmortalidad, entonces Hera se lo quita de encima derramando la leche y creando la Vía Láctea.

\chapter{Clase 4}\label{ch:clase4}
Belerofonte es un héroe con relevancia en la época clásica, destaca por su montura Pegaso, aunque en la actualidad no. Pegaso se ha representado en la época Helenística.

El nombre de Belerofonte viene de cuando mato a su hermano, se oculta y el rey le refugio, la reina se enfada de celos y lo manda matar por el asesinato de su propio hermano.

También aparece como que Belerofonte recibe una carta que dice que debe matar al que le haya entregado la carta, típico leitmotiv.

Yóbates le encargan matar a una quimera, cuerpo de león y cola de serpiente. La consigue matar tras pedir ayuda a los dioses, Atenea le ayuda con unas bridas especiales a capturar a Pegaso, con el que crea un fuerte vínculo, este le ayuda a matar al monstruo. Tras esto le nombran rey. 

En el momento culminante de su gloria, trata de subir con Pegaso al Olimpo, pero le mandan un tábano que le pica y tira, muriendo o vagando por la tierra maldito.

\section{Hera - Juno}
\begin{itemize}
    \item Soberana del Olimpo.
    \item Mujer y hermana de Zeus.
    \item No madre de dioses (provincia compartida).
    \item Diosa del hogar matrimonio. Es el símbolo del matrimonio romano, en el que sufre idas y venidas, soportando al marido que está por ahí.
    \item Hera Ziguía, con faceta de soberana, matrona y protectora. Es una figura secundaria.
    \item Faceta mitológica: celos maritales.
    \item Diosa antigua. Algunos de los templos más antiguos estand dedicados a ella.
    \item Sin evolución cronológica clásica.
    \item Se le sacrificaban vacas y a Zeus toros.
\end{itemize}
\pagebreak

\section{Heracles ''el que place a Hera''}
Los trabajos de Hércules (Hera le hace enloquecer, matando a su familia, hijos o mujer dependiendo de la versión) para redimirse, típico de los héroes, los dioses le dan proezas:
\begin{itemize}
    \item Errores
    \item Purga de sus pecados (asesinato de hijos)
    \item Trabaja para Euristeo, rey de Tirinto, que le asigna 12 (número simbólico) trabajo imposibles.
    \item Solo o con ayuda
    \begin{itemize}
        \item Divina (variada)
        \item Heroica (Yolao)
    \end{itemize}
    \item Muerte trágica
    \item Sublimación final
    \begin{itemize}
        \item Cuerpo mortal
        \item Ascenso al Olimpo y reconciliación
    \end{itemize}
\end{itemize}
El primero es la caza del león de Nemea, que lo ahoga por su piel es permeable a las armas por eso la usara y se hará símbolo de él.

Muerte: Heracles se enamora de Yanira y se queda con ella para vivir, tras matar a Neso (un centauro malo, opuesto a Quirón) y regalarle una capa llena de sangre de este. Al tiempo Deianira su esposa, de celos le regala una capa envenenada a Yanira que se la entrega a Heracles, quien al ponérsela empieza a quemarse. Sufre tal dolor que pide a los dioses que acaben con su vida, se lo conceden y su parte de dios sube al Olimpo y se reconcilia con su madre.

Antropofagia y el fuero eran signos de la barbarie.

Apolo y Heracles luchan por el trípode délfico, de Delfos, finalmente gana Apolo, pero él resuelve la duda de como librarse de la maldición de la muerte de Ífito.
\begin{itemize}
    \item La sacerdotisa se sentaba sobre él y entraba en trance.
\end{itemize}

Héroe Dorio Civilizador (Hércules es el dios civilizador por antonomasia, cuanto más bajo caes más puedes subir)
\begin{itemize}
    \item Mitos secundarios
    \item Heráclidas
    \begin{itemize}
        \item Descendencia
    \end{itemize}
    \item La civilización alcanza  los límites del mundo conocido
\end{itemize}

Hércules pide a Atlas que le traiga unas manzanas de sus hijas, mientras él aguantaría el mundo, pero no vuelve... Entonces Heracles le dice que le sujete la esfera un momento y en este momento se zafa de sujetarlo más y le roba las manzanas. 

Es estrecho de Gibraltar, las columnas de Hércules.

La dinastía de Esparta decían ser descendientes de Heracles.

\section{La Apoteosis (sublimación)}
\begin{itemize}
    \item Referencia de sublimación, paso a ser dioses o diosificados, culto al líder y asciende a los cielos para ser dioses. Ascienden como ceniza tras su muerte o llevados por Marte como los del origen de Roma
    \item Heracles - Hércules
    \item Reyes helenísticos
    \item Emperadores romanos
    \item Neoclasicismo: Apoteosis de héroes nacionales.
\end{itemize}

Culto latino

Heracles y Caco (significa el mal), se le erige el Templo Hércules Victor (vencedor)

Foro Boario

Dios popular

Templo púnicos (Melkart)

Cómodo se hizo representar como Hércules a finales de siglo II d.C.

Los trabajos de Hércules ya en el mundo latino se difundieron por todo el Mediterráneo. Como un dios fuerte y heroico. Con el cristianismo baja su popularidad.

\section{Hércules en Hispania: Gerión, Pyrenne y Tubal}
\begin{itemize}
    \item Monstruo Gerión.
    \item Ninfa Pyrenne. Por la pena de Heracles tras la muerte de Pyrenne empieza a amontonar piedra y se da origen a los Pirineos.
    \item Personaje bíblico menor, Túbal, nieto de Noé. Se repartieron el mundo y se encarga de poblar la península ibérica, dando entre otros a los vascos.
\end{itemize}

\chapter{Clase 5}\label{ch:clase5}
Tema clase: Desde los dioses hasta los hombres. Las edades del hombre o de la humanidad.

Los humanos son hijos de la Tierra, Gea, u originarios de diferentes regiones.

\section{Otros orígenes las edades del hombre}
\begin{enumerate}
    \item Hesíodo frente a Ovidio
    \begin{enumerate}
        \item \textbf{Oro - Cronos.} Sin sufrimiento, se alimentaban de lo que la naturaleza proporciona
        \item \textbf{Plata - olímpicos.} 100 años de juventud. Soberbia, condena al Hades, no rendían culto a los dioses. Son castigados con un diluvio universal. Pasan a tener que trabajar y envejecen rápidamente.
        \item \textbf{Bronce - Zeus.} Guerreros y cultura de bronce (como material), fuertes y valientes para la batalla. Se alimentan con carne y no solo cereales. Autodestrucción. Hades sin memoria.
        \item \textbf{Héroes - héroes del periodo mítico.} Desaparecidos o en los márgenes. Las hazañas a diferencia que el bronce son recordadas por los hombres. No tiene nombre de metal porque no es reconocida por Ovidio.
        \item \textbf{Hierro - actual.} Pecadora y codiciosa, acechada por los males para ellos, el medioambiente y los dioses. Material más resistente.
    \end{enumerate}
    \item Fin del mundo: cuando el bien, juramento y justicia no tengan valor.
\end{enumerate}

Uno de los mitemas es la degeneración de la sociedad, siempre la sociedad previa era mejor, según se van alejando de su vida en el Edén. Ciclos de creación y destrucción. Los mexicas o aztecas usaron este tema como eje de su teología política y los sacrificios se hacían para evitar el fin.

\section{Poseidón / Neptuno}
\begin{itemize}
    \item Soberano de los mares
    \item Agua dulce
    \item Terremotos y caballos
    \item Tridente y carro de hipocampos
\end{itemize}

Muchos de sus descendientes son terribles.

Poseidón tenía los juegos del istmo de Corinto.

Dio lugar a ninfas, pueblan la naturaleza como diosas menores. Supervisoras de lugares sagrados, en Galicia hay muchas representaciones. Ninfeos, ...

No tuvo muchos lugares de culto, además de Corintia. Atenea confronta a Poseidón para hacerse con Atenas.

Las ciudades tienen un dios patrones que tiene un lugar preferente.

Atenas es el centro de la Grecia clásica, donde se conectaba la ciudad con el puerto.

En la Atenas clásica era el centro cultural y político del momento, donde se escriben importantes obras como la Ilíada. Tenía la acrópolis, el ágora, ... Tenían una gran estatua de Atenea.

\section{Teseo: Héroe de Atenas}
\begin{itemize}
    \item Bastardo del rey Egeo. Cuando llega su adolescencia es llevado por su madre a recoger bajo una pesada piedra su armadura y arma, enterrada por su padre para cuando fuese digno.
    \item Fuerte e inteligente (ingenio), no se dejó engañar por el asesino de Corinto. 
    \item Héroe jonio
    \item Mitos históricos
    \item Civilizo el camino a Atenas. Venció a múltiples bandidos con nombre, el héroe civiliza las ciudades para que los hombres las puedan poblar.
\end{itemize}

\section{Creta y Monte Ida}
\begin{itemize}
    \item Ida y yacimiento de Zominthos. En el monte Ida hay una cueva en la que se han encontrado ofrendas a Zeus, además era un centro de peregrinaje al lugar donde creció Zeus
    \item Juventud de Zeus
    \item Refugio de Europa
    \item Casada con el rey y 3 reglados: perro, jabalina y automata (el automata Talos, parece se una figura importante, era defensor y tenía su punto deben en el tobillo)
\end{itemize}

\section{Hades}
\begin{itemize}
    \item Hijos de Europa: Radamantis, Sarpedon y Minos.
    \item Regalo de Poseidón: el toro blanco. Se hizo con parte del mediterráneo, pero por su soberbia y no querer sacrificarlo es castigada su mujer.
    \item Relación con la cultura minoica
\end{itemize}

\section{Un castigo divino: Pasífae y el Toro}
\begin{itemize}
    \item Pasífae es asistida por Dédalo para entrar en la vaca de madera y acercarse al Toro de Poseidón, teniendo relaciones con el, dando como resultado el minotauro. Creció rápidamente y demostró su naturaleza salvaje, alimentándose de personas y matando ganado.
\end{itemize}

\section{El laberinto}
\begin{itemize}
    \item Labrys: hacha minoica y objeto ritual.
    \item Palacio y símbolo de identidad posterior en Creta
    \item Dédalo le construye este palacio en Labrys, de ahí laberinto.
    \item El rey Minos cada ciertos años mandaba a 14 jóvenes como alimento, no lo querían matar de hambre.
    \item Teseo es el encargado de vencer al Minotauro, acabando con los sacrificios y convirtiéndose en rey. Previamente, es retado por el rey a recuperar un anillo tirado al centro del océano, llevándose además una corona.
    \item Ariadna se enamora de Teseo  y le cuanta el secreto para atravesar el laberinto. Hilo de Ariadna. Teseo era capaz de vencerlo, pero no de salir.
    \item Egeo (padre Teseo) se queda compungido al ver las velas negras del barco que volvía y se tira al mar. Dado que se suponía que llevaría si gana velas blancas y de la tristeza se tira al mar, así surge el mar Egeo.
    \item En algunas versiones no muere el minotauro, sino que lo ata y lo llevara celebrando su victoria, lo que daría lugar a ciertos cultos dedicados al toro.
\end{itemize}

Ícaro es encerrado con su padre Dédalo, dado que construyeron el laberinto y sabían salir de este y de la isla. En esta torre construyen un par de alas de cera cada uno, ambos huyen, pero advertido por su padre se acerca demasiado al calor de Helios se le derriten y caen. 

Minos estuvo buscando al inventor huido, Dédalo, escondido en Griganto. Pero con un truco, hilar una concha, imposible para el resto de mortales, consiguió que por su orgullo saliera y lo lograse (usa una hormiga que tira de ella). 

En un baño Minos es ahogado, siguiendo Dédalo con su vida, pero termina la época.

\section{Minos}
\begin{itemize}
    \item Muerta por Artemisa (Homero)
    \item Abandono en Naxos (después)
    \item Rescate de Dionisio volviendo de la India (lugar lejano y místico, desde donde los dioses traen nuevos conocimientos)
    \item Ariadna se convierte en su esposa
\end{itemize}

\section{Otras aventuras de Teseo }
(hay muchas historias y variaciones, muchas inconexas)
\begin{itemize}
    \item Con las amazonas con Hércules, rapta a Hipólita y nace Hipólito. Esta será maltratada y posteriormente abandonada.
    \item Pirítoo e Hipodamia y la Centauromaquia
    \item Teseo se casa con Helena de Esparta que más tarde seria Helena de Troya
    \item Descenso a los infiernos, donde Hades les tiende una trampa y los deja pegados a un banco por artes mágicas (cinto mágico o brazalete, según la versión) como castigo por toda la eternidad. Heracles en su duodécimo baja al Hades y se los encuentra, como son amigos logra liberarlos con su fuerza. Con esa fuerza tiembla la Tierra y abandonan a Pirítoo en el infierno.
    \item La madre de Teseo es vendida como esclava y hacen huir a sus hijos con Fedra. Poniendo en el trono a un enemigo de Atenas.
    \item Teseo en la cúspide de su gloria pierde todo por intentar meterse donde no le llaman, con los dioses.
\end{itemize}


\chapter{Clase 6}\label{ch:clase6}
Siguiendo con Teseo\dots

El reconocimiento de los atenienses y otras partes de la alianza ateniense es lo que dimensiona ese mito local y lo engloba junto al resto del mito de Teseo, haciendo que se sepa históricamente, está puesto en los libros de historia griegos que los huesos fueron puestos en el Teseion.

\section{El Teseion}
\begin{itemize}
	\item Culto heroico
	\item Culto cívico ateniense a Teseo
	\item Sustituto de ancestros
	\item Huesos traídos desde Esciro en el año 475 a.C.
	\item Es el lugar donde se concentra toda la moral cívica mezclada con religión y tradición local.
	\item Es el gran templo dedicado al culto a Teseo que se sitúa en la acrópolis ateniense, aunque en escritos de la época bizantina mucho después no se sabe específicamente donde estaba.
\end{itemize}

Teseo de Creta se convirtió en el elemento por excelencia para justificar la singularidad ateniense, como también ocurrió a los romanos con Rómulo y Remo para hablar de su característica especial y porque dominan el mundo. Los estadounidenses hablando del american way of life donde se crearan una serie de herramientas de libertad democrática tratadas casi como elemento religiosos, casi míticos, de alegoría a una moral más que a uno fundacional de una ciudad.

Igual que ocurriría en el mundo cristiano, donde algunas ciudades se relacionan muy fuertemente con su santo patrón, donde la gente tiene una relación sentimental con este santo patrón.

\section{Perseo}
\begin{itemize}
	\item Zeus padre, Acrisio abuelo
	\item Destino profetizado
	\item Niño abandonado
	\item Encargo mortal
	\item Dones de los dioses
\end{itemize}

El rey de Argos, que moriría a manos de su propio nieto, para evitarlo adquirió e hizo que su hija Dánae fuera encerrada en una mazmorra de bronce o una torre alta de bronce (según las versiones) para impedir que tuviera ningún tipo de trato con el exterior, especialmente con ningún varón que la dejara embarazada.

Pero el rey de los dioses Zeus sí que la vio y sí que pudo acceder, se transformó en lluvia dorada y cayó a través de la ventana sobre Dánae preñándola. De este embarazo milagroso nace Prometeo.

El rey Acrisio que sabe que es un semidiós, tuvo miedo de matarlo directamente, por lo que lo coloca en una caja de madera a él y a su madre y los echa al mar, donde llegan a la isla de Sefilos y son regidos por un amable rey. Perseo creció en el exilio como un príncipe bastardo, pero el mismo rey se enamoró de Dánae y ante el rechazo de esta que solo quería criar a su hijo decidió deshacerse de Perseo al ser un obstáculo para conseguir su amor, por lo que le pide una tarea imposible que le traiga la cabeza de medusa y ante la desesperación del héroe sabiendo que es casi imposible este reza los dioses.

Atenea y Hermes bajan de los cielos para ayudar al joven semidiós, le dan un yelmo de Hades que lo hacía invisible, Hermes le prestó sus sandalias aladas, una espada especial muy afilada irrompible para cortar la cabeza de medusa y por último le entregó una bolsa donde guardar con seguridad la cabeza del monstruo.

Atenea le aconseja volar al oeste con estos regalos, hacia la oscuridad lejana, donde podrá encontrar a tres brujas, las Greas (no confundir con las moiras) que comparten un mismo ojo.

Perseo para convencerlas lo que hace es acercarse a ellas, hablarles amablemente y les roba el ojo, por tanto, les obliga a cambio del ojo mostrarle el camino para encontrar a la gorgona medusa. Cuando llegó a la tierra de las gorgonas también eran un grupo de 3 y solo una de ellas será mortal. Decide matarla esperando a que estén dormidas, pero Atenea aparece de nuevo en ese momento para evitar que las despierte y le dice que espere que se acerque sigilosamente y le corte la cabeza a Medusa mirando su cara a través de la imagen que se refleja en su escudo para evitar que la cara de medusa lo volviese en piedra.

La cabeza tenía víboras como pelo, la corta y la mete en la bolsa alejándose, pero de las gotas de sangre que sale en el cuello medusa surge Pegaso, el caballo alado.

En su vuelta pasa por Etiopia y ve a un dragón enviado por Poseidón que trata de devorar a la princesa Andrómeda, que había sido ofrecida a este por haber ofendido a Poseidón diciendo que era más bella que las nereidas. Perseo llega volando, mata al monstruo (con la espada de los dioses o la cabeza de medusa, según la versión) y rescata a la princesa. Entonces queda prendado de ella y decide casarse con Andrómeda.

En la boda tratan de recuperar a la princesa, pero usa la cabeza de medusa y convierte a la mitad en piedra, por la situación se ve obligado a huir.

En Argos como héroe, en unos juegos en su honor, compitiendo, tira un disco (el que más lejos) y cumple la profecía, matando de un golpe en la cabeza a Crisio sin querer. Se le ofrece el trono, pero por vergüenza de este asesinato no rechaza, en la época aunque no fuera intencional es culpable.

Se irá y funda Micenas, fundamental en los ciclos troyanos.

\section{Atenea Minerva}
\begin{itemize}
	\item Hija de Zeus y variantes
	\item Diosa de la sabiduría y olivo
	\item Sabia consejera
	\item Portadora de la Egida
\end{itemize}

Atenea de Atenas visitó el lugar y de allí trajo ideas como la justicia. De hecho, también se dice que esta diosa trajo sus ropajes extraños de libia, por eso aparece con un aspecto de pieles de cabra adornadas con serpientes, que será la vestimenta original de las mujeres antiguas de roma.

En el mito famoso el más famoso de todos es hija de Zeus y Metis, la titán, la primera mujer de Zeus que había sido profetizada que engendraría un hijo que castigaría a Zeus y lo vencería. Y como Cronos había derrotado a su padre, Urano, alarmado por esto, lo que hace es devorar a su esposa Metis. Por esta razón Zeus recibe la razón de la razón justiciera, que es lo que significa Metis y ella le aconseja desde sus tripas. A la vez empieza a experimentar unas migrañas terribles y decide apartarse del olimpo para intentar descansar y relajarse, pero estas migrañas no se van. Hermes entonces acude en su ayuda y adivina qué es lo que le pasa, Metis está embarazada y la hija ha nacido y está intentando salir. Entonces llamada a Hefesto y junto con él abren una brecha en el cráneo de Zeus y de allí sale ya lista para la batalla y armada la diosa Atenea.

El dios se come a en este caso Zeus no mata Metis y acabó con embarazo, sino que se la come pensando lo que acabará con el embarazo, es una parte simbólica de absorber la fuerza del vencido.

Los símbolos son el búho mochuelo de la sabiduría, el árbol del olivo por el mito de Atenas y la serpiente de la gorgona.

Se ha relacionado con Minerva, pero Minerva realmente era originalmente una diosa itálica distinta, una mecánica común que hay júpiter no es el mismo que Zeus. El conjunto de iconografía de los griegos, pero también este caso el mitológico y el religioso, de esta forma encontramos que Atenea y Minerva acaban siendo progresión de la misma diosa, se vuelven equiparables.

\section{Troya}
\begin{itemize}
	\item Combaten griegos y troyanos
	\item Caballo de Troya
	\item Muere Héctor a manos de Aquiles
\end{itemize}

Es el elemento fundacional de la cultura griega y por esa razón una de las semillas principales de la cultura occidental, no solo europea sino también estadounidense y de Latinoamérica directa e indirectamente.

En la mitología, en este caso, Troya engloba numerosos episodios de los cuales uno, el más famoso, La Guerra de Troya, es el nombre general y la Ilíada de Homero es la obra particular que no es lo mismo. La Ilíada es la obra principal para los griegos, es una obra cultural de referencia

\section{Tetis y Peleo}
Se obliga a esta diosa menor a casarse con un mortal para evitar que nazca un dios peligroso y el mortal seleccionado no fue ni más ni menos que peleó el rey de los mil millones en tesalia.

La boda
\begin{itemize}
	\item Tetis y dioses
	\item Peleo tesalio y reyes
	\item Excusión de Eris
	\item Juicio de Paris
\end{itemize}
La diosa no quería casarse y huyó constantemente y Tetis se transformaba en animales para escapar igual que su padre Nereo.

Un matrimonio no correspondido del que nacerá Aquiles y en este momento cuando se consolida la unión a un momento antes del nacimiento de Aquiles se celebra la boda.

Lanza una manzana con la inscripción para la más bella, donde las diosas Minerva (Atenea), Hera y Afrodita (Venus) se quedan solas en la disputa por la manzana, en esta situación se recurre a París un joven y apuesto príncipe troyano que aparece siempre como pastor y Hermes le lleva la manzana para elegir la más bella. Cada una de las diosas le ofrece algo, Hera le ofrece el dominio sobre los hombres, Minerva le ofrece la sabiduría mayor que a cualquier hombre y Afrodita le ofrece el amor de la mujer más bella.

Afrodita le dará el amor de Helena a París, aquí también donde pasamos a hablar del rapto de Helena, aunque hay diferentes versiones algunos dicen que no fue llevada por la fuerza, sino que se dejó enamorar o quedó prendidamente enamorada por París. Helena, que era hija de Zeus y de le da, por tanto, era también semidiosa.

El rey de Micenas (ciudad que ha sido fundada por Perseo) y señor de todos los reyes griegos son llamados a la guerra para recuperar a Helena y para evitar el agravio de los troyanos.

Aquiles es para ese momento ya había empezado a crecer como uno de los héroes más importantes del mundo griego, de hecho se dice que Peleo, su padre, lo entregó a los centauros que vivían en las fronteras de tesalia y fue educado por ni más ni menos que Quirón, el centauro sabio que formaba a los héroes.

Le dan a elegir entre una vida larga o una vida gloriosa, y decide la muerte gloriosa para alcanzar la inmortalidad a través de su fama.

Lo que hace la Ilíada es mencionar la guerra de Troya a mitad de la guerra, es un conflicto gigantesco, lo que buscaba este relato era acabar con la hiperabundancia de humanos en el mundo por todo lo que había ocurrido desde el matrimonio de Peleo y Tetis, la manzana de la discordia y el enfrentamiento de las tres diosas. Era una artimaña de Zeus para que los humanos se enfrentarán entre sí y los dos grandes grupos de población en mundo griego y Asia se enfrentasen y redujesen la población y su fuerza en gran medida para calmar la hiperpoblación del mundo.

La cólera de Aquiles que es el dinamizador de la mayoría de relatos, este criado por Quirón se enfada y consigue grandes proezas, pero también grandes daños que es lo que ellos pretenden. La cólera sin control que prácticamente destruye el mundo narra la guerra de Troya, todo cumpliendo la voluntad de Zeus.

Los troyanos se acercan a la victoria definitiva y acorralando contra la playa a los griegos, entre ellos está el gran amigo o amante (una constante en el mundo griego, relación entre la parte poderosa y la parte querida más débil) Patroclo decide unirse a las fuerzas en algunas explicaciones contra y otras con los deseos de Aquiles toma sus armas y sin la participación de Aquiles intenta el último defender el campamento griego y detiene el avance troyano, pero a costa de su vida, ya que se enfrentará con el gran héroe troyano Héctor y recibirá muerte. En este momento es cuando la cólera de Aquiles alcanza una sobredimensión, una cólera divina, Aquiles pide ayuda a su madre y a los dioses, estos le regalan una nueva armadura y lleno de furia vence a los héroes de distintas ciudades de Troya.

Acorrala contra las murallas de Troya a los troyanos y empieza a buscar al asesino de su querido Patroclo, Héctor, es aquí el destino definitivo de uno de los grandes héroes y la cúspide del relato de la Ilíada.

Héctor es el héroe cívico por antonomasia, luchado por su ciudad, por su familia, por sus gentes y además se le muestra con paciencia. Todo se sale de su sitio cuando finalmente se tiene que enfrentar o piensa que está enfrentando Aquiles y mata a Patroclo, entonces Aquiles lo perseguirá. Héctor decide evitar el conflicto porque sabe que no está bajo una rabia normal, sino que está bajo una rabia divina e intenta huir, la Ilíada cuenta como dan vueltas a la ciudad de Troya, esperando a que perdiera su furia inicial para combatir.

Engañado por Atena que se hace pasar por su hermano que le pasa las lanzas, es derrotado por Aquiles.

El héroe troyano lacera sus tendones, los ata a su carro y desde los tobillos lo arrastran desde la ciudad de Troya extramuros durante doce días expuesto el cadáver al sol y los animales.

En una pira al que acuden todos los héroes a la vez en una tregua sagrada donde observando como el cuerpo de Héctor se va deshaciendo pasto de las llamas, acaba el relato épico de la Ilíada.

\chapter{Clase 7}\label{ch:clase7}

\section{Otros relatos del ciclo troyano}
Tras la muerte de Héctor se trae a los bárbaros y las amazonas (hasta que derrotan a su líder) para combatir a los troyanos.

\section{El destino de los héroes}
El fin de la Ilíada es la muerte de Aquiles de un flechazo en el tendón, el talón de Aquiles. Lo hace Paris tras la muerte de su hermano, mediante varias flechas.

No se sabe en qué momento se redactan sus muertes, así como la de Paris u otros personajes.

Se dice que es débil en el talón porque no lo sumergió cuando lo hicieron inmortal o Tetis lo alimenta con ambrosía.

\section{El final de La Guerra de Troya}
Ulises convence a los reyes de que hagan una treta para hacerse con la victoria, ya que no avanzaban en la guerra, simulan una huida a unas islas de Turquía, pero unos pocos hombres se esconden en una ofrenda a los troyanos. Por la noche salen y masacran a los troyanos.

\section{Jason y los Argonautas}
\begin{itemize}
	\item Ciclo pre-troyano. Estos héroes viajan por un mundo que se va abriendo lugar y van descubriendo elementos de los que se hablara más adelante en la Odisea y los ciclos.
	\item Nefele y Atamante: el vellocino de otros. Toisón de oro (tomado por la corona francesa). Es un carnero alado que rescata a los herederos al trono que van a ser sacrificados para salvar la hambruna, se dice que es enviado por los dioses. Por el camino en el mar negro se cae y muere uno de ellos y el otro llega con éxito.
	\item Colquida (actual Georgia) y Eetes
	\item Jason de Yolco, Tesalia. Es criado por Quirón y es probado por Hera (vieja que cruza un río) que le ayudará.
	\item Hermanastro Pelias y se le pone una tarea heroica, traer el vellocino, de esta manera se gana el puesto al trono (aunque era legítimo). Se le pone por la profecía que tenía el rey de que un forastero con una sola sandalia le destronara.
	\item Argonautas y viaje mítico. Atenea ayuda a construir el barco Argo con un tablón, con el que emprenderá el viaje, como tripulación se unen los argonautas, una serie de héroes y personajes relevantes (Padre de Aquiles, Heracles, hijo de Hermes, ...). Es un viaje con muchas etapas, que se cuentan como historias y mitos relacionados con las etapas.
	\item Algunas veces son recibidos como aliados y otra como enemigos, como en Micos, el hijo de Poseidón que era boxeador (los hijos de Poseidón son monstruos o agresivos)
	\item Se le entregaría el vellocino si Jason hace 3 tareas, pero cae en depresión y no la completa, por lo que con ayuda de los dioses enamora a la princesa (Medeas) quien le ayuda y logra sus tareas.
\end{itemize}

Hay una representación en el cine de uno de estos viajes míticos.

La historia de Medea: la mala mujer. La naturaleza más oscura del amor, fertilidad, la brujería y los sentimientos derivados: engaño, celos y resentimiento. Mata a la mujer de Jason con un vestido maldito, que se pega a la piel y la quema. Además, mata a sus hijos.

\section{Afrodita (Venus)}
\begin{itemize}
	\item Hija de la espuma. La más hermosa elegida por París.
	\item Diosa del amor y erotismo. Es una personificación.
	\item Popular en la mitología y alegoría, más que en los cultos.
	\item Eros es hijo suyo junto a Hades.
	\item Esposa de Hefesto, que al enterarse de su aventura con Hades los atrapa y expone en la cama.
	\item Guerra de Troya y otros mitos
	\item Festival Afrodisias, tiene una representación más carnal en algunos lugares y en otro más sentimental. Aparecían en varias ciudades griegas, pero eran especialmente relevantes en la isla de Chipre, lugar de origen de la diosa según la mitología.

	      Era muy popular en especial entre las prostitutas, que la consideraban su patrona.
\end{itemize}

Tiene relación con dos mortales, ...  y Adonis del que se enamora de verdad. A la muerte de Adonis en una cacería por el ataque de un animal, Afrodita corre descalza a recoger su moribundo cuerpo y de las heridas de sus pies nace el rojo de las rosas o las anémonas.

\section{Odiseo (Ulises) y la Odisea}
Se caracteriza por su sabiduría e ingenio. Criado por Quirón y casado con Penélope.
\begin{itemize}
	\item Telemaquia, nos cuenta la historia de los héroes que se van repartiendo por el mundo.
	\item Viajes de Odiseo.
	\item La Odisea comienza en mitad del viaje, in media res.
	\item La venganza de Odiseo
\end{itemize}

El viaje es de regreso a Itaca, su hogar.

Odiseo llega a la isla de los cíclopes donde en una cueva encuentran alimento y se dan un banquete, entonces llega Polifemo, hijo de Zeus y se los empieza en comer en venganza. Ulises con sus hombres entonces lo emborrachan y le ciegan con una lanza. Por lo que pueden escapar cuando saca a pastar a sus ovejas. Zeus le maldice y deberá navegar otros 10 años.

En la isla de Circe, bruja semidiosa, les convierte en cerdos, pero gracias a un antídoto de Hermes se vuelve humano y la obligan a devolverlos a su forma. Llega a vivir allí un año.

Por el viaje encuentran sirenas mitologías y son atados a los mástiles, aquellos que no reman para que no se tiren al mar y sean devorados por estas bestias.

Ulises desciende al inframundo en busca de Tiresias para que le aconseje e ilumine acerca de su regreso a Ítaca. Nada más llegar, ve acercarse una multitud, la de los que no son personas ni tienen rostro, no son visibles, no son nada. Entre ellas distingue el espectro de Aquiles, al que da de beber sangre para devolverle algo de vitalidad y pueda expresarse. El héroe aqueo le dice que los muertos están privados de sentidos y que son las imágenes de los hombres que ya fallecieron. Añade que preferiría ser el último servidor del hombre más pobre del mundo, pero vivo bajo la luz del sol, que ser el rey de ese mundo de tinieblas que es el Hades.

Tras 20 años regresa a Itaca, pero nadie le reconoce más que su perro. Cuando ve a Penélope, esta no le reconoce y le dice que había anunciado que quien consiguiera armar el arco se casaría con ella. Penélope sabia que solo sería capaz de hacerlo su marido, Ulises, porque no quería casarse con otro hombre. Ayudado por sus hijos y sirvientes, acaba con los pretendientes.

Es uno de los pocos relatos en los que el héroe acaba bien.

Esta obra tiene mucha influencia en la cultura.

\chapter{Clase 8}\label{ch:clase8}
\section{Rómulo y Remo}
Es uno de los clásicos mitos fundacionales, que tiene origen con Alba longa fundada por Eneas. En Alba longa reinaron dos hermanos, para que solo reinara otros le expulsa.

Son los hijos gemelos de Rea, Silvia y Marte (descendientes directos de Eneas) que fueron rescatados en las orillas del río Tiber (zona pantanosa) por una loba (loba capitolina/luperca), que se los llevo, como esta había perdido sus lobeznos, los acoge y los amamanta como a sus propios hijos. Se dice que originalmente podrían haber sido por una prostituta dado que la palabra significa lo mismo que loba.

Los orígenes de Romas son el fango y la loba.

Un pastor, Lumitor, les explica su origen de alta cuna y estos volverán a Alba longa para vengarse junto a un conjunto de pastores. Una vez con esta venganza se van para fundar su propia ciudad, su localización seria donde la loba los rescata. Como no se ponen de acuerdo, el que viera más buitres decidiría el nombre.

En una pelea entre los hermanos acaba muriendo Remo, que fue enterrado en el punto de origen de la ciudad de Roma. Fue ayudado por los patricios.

El rapto de las Sabinas, como había pocas mujeres para poblar roma, cuando los sabinos vinieron a la ciudad raptaron a sus mujeres. Con el tiempo las mujeres raptadas se adaptan y se les da ciudadanía, como uno más, teniendo a sus propios hijos y maridos. Hubo una gran lucha entre los romanos y los sabinos, pero fue detenida por estas mujeres dado que ganara quien ganara perderían (morirían sus maridos o sus familiares).

Cuando Rómulo muere comienza la monarquía, a partir de aquí hay varias versiones, alguna de ellas dice que es matado por los patricios y otra que asciende a los dioses.

\section{Marte - Ares}
\begin{itemize}
	\item Dios de la Guerra
	\item Papel menor en Grecia, naturaleza distinta en Roma
	\item Relación con Venus (Afrodita)
	\item Unido a Bellona (diosa de la guerra) en Roma, en la mitología griega no tiene una mujer asociada
	\item Hijos varios, destacan Fobos y Dimos
	\item Se le considera muy importante por su vigor bélico, también como dios que ayuda a Rómulo a ascender
	\item Sus símbolos son la lanza y el fuego
\end{itemize}

\section{Los dioses de la Luz}
\subsection{Apolo}
Es joven y representaba la belleza, masculinidad y es la luz sin ser Helios. Representa también la sanación, asociándosele peanos (un peán). Se le asocia como animal el ratón y como instrumento la lira, que le es entregada por un dios.
Hijo de Zeus y Leto (titanide), de ahí nace Apolo y Artemisa. Fue un parto muy duro al no contar con la presencia de Hera y la representación del parto.

Dafne significa laurel, que se usaba como símbolo de la victoria.

\subsection{Delfos}
\begin{itemize}
	\item Es el templo más importante de la griega occidental
	\item Templo de Apolo en Delfos, sede del Oráculo
	\item Pitonisa, daba oráculos sobre el trípode. Las pitonisas mascaban laurel y daban repuesta a las cuestiones que les plantean
	\item Delfos era el ombligo religioso y geográfico del mundo griego (ónfalos), Apolo lanza dos buitres opuestos y se cruzan ahí
\end{itemize}

Zeus y Apolo son los principales dioses oraculares.

\subsection{Apolo y las Artes}
Las musas son divinidades inspiradoras de las artes: cada una de ellas está relacionada con ramas artísticas y del conocimiento. Son hijas de Zeus y de Mnemósine, compañeras del séquito de Apolo.
\begin{enumerate}
	\item Calíope - elocuencia y poesía explica
	\item Clío - historia
	\item Erató - poesía
	\item Euterpe - música
	\item Melpomene - tragedia
	\item Polimnia - himno
	\item Talía - comedia
	\item Terpsicore - danza
	\item Urania - astronomía, didáctica y las ciencias (ciencias exactas)
\end{enumerate}

\subsection{Helios (Sol)}
Es un titán más que un dios, es hijo de Hiperion y Eurifae (Tea). Es el grecolatino Sol. Se caracteriza por su carro, cuadriga solar, que le permitía moverlo por el cielo.

Suele aparecer con un limbo en la cabeza o luz, que será utilizada también por el cristianismo. Ese halo, además de estar en la divinidad cristiana, deriva en la corona que se utiliza en la monarquía.

Da origen a muchos personajes que hemos visto, como a Circe o Medeas, que les da poder para poder, en caso de Circe para convertir en cerdos. También Pasífae, la madre del minotauro de Minos y Serena/Luna. Otro destacado son las Horas (4 para estaciones, otras veces por 12 por las horas de día y de noche), aunque muchas veces son atribuidas a Zeus, así como las Tres Gracias.

\section{Mitraísmo}
Culto que tiene su propia mitología copiada del culto Persa. Narran en unos relatos una alianza entre Mitra y Sol para derrotar al toro cósmico. Hay diferentes niveles de pertenencia a este culto, desde iniciado hasta páter.

\section{Faetón}
\begin{itemize}
	\item Mito trágico del recorrido solar
	\item Evolución del mito y su estructura
	\item Mito moralista (fábula)
\end{itemize}
Joven, bello, fuerte y apuesto. Se burlaban de él diciendo que no era hijo de Helios, pero su madre le dice que es cierto, que vaya hasta el final mundo para verle. Faetón le pide una prueba de paternidad, Helios dice que le dará lo que sea, entonces le deja montar en la cuadriga. Se arrepiente, pero como se lo promete le deja advirtiéndole, ni los dioses querían conducirla.

Faetón lo monta y pierde el descontrol, subiendo demasiado revolviendo las galaxias, entonces baja demasiado y quema la superficie de la Tierra, secando el suelo, ríos y quemando a los etiopes (volviéndose negros). Por este descontrol, los dioses le tiran un rayo para sacarle de la cuadriga y muere al caer. Entonces es enterrado por sus hermanas, la heliades  que se convierten en álamos.

Al día siguiente, Helios no salió y hubo un día de oscuridad (eclipse).

\end{document}